\section{The no-fold branch and the infinite tail}\label{sec:terminal}
Fix the phase interval $[a,d]$ containing $W$. Alongside $P_t$, its
certificate gives a positive summable sequence $(u_t)_{t\ge1}$ such that
\begin{equation}\label{eq:supersequence}
 u_1\ge M_1(a),\qquad u_2\ge M_2(a),\qquad
                  \sum_{s\ge t-2}u_s\le P_tu_t\quad(t\ge3).
\end{equation}
The initial terms are rational. The continuation is the positive recurrence
\begin{equation}\label{eq:tail-recurrence}
 u_{t+2}=(1-r)u_{t+1}+r(1-r)u_t\quad(t\ge15),
 \qquad r=\frac{88331750161}{10^{11}}.
\end{equation}
Since $0<r<1$, the recurrence is summable. If $\beta=r(1-r)$ and
$Q=1/[r^2(1-r)]$, then
\begin{equation}\label{eq:tail-sum}
       \sum_{s\ge15}u_s=Q\bigl(\beta u_{15}+(1-r)u_{16}\bigr).
\end{equation}
This follows by summing the recurrence. Its positive characteristic root
is less than $1$, or equivalently the same conclusion follows by the
geometric majorant used in the formal tail proof. The certificates check
the finite prefix of~\eqref{eq:supersequence} and the recurrence conditions
that extend it to every $t$; in particular $Q<P_{16}$.

\begin{lemma}[No-fold mass bound]\label{lem:no-fold-mass}
If no active stage with positive core mass satisfies~\eqref{eq:fold-test},
then
\[
                    S_i\le n u_{m_0-i}\quad(0\le i<m_0).
\]
In particular the cores retained before the terminal stop have total mass
at most $n\sum_{s=3}^{m_0}u_s$.
\end{lemma}
\begin{proof}
If the conclusion fails, take the maximum
$R=\max_i S_i/(nu_{m_0-i})>1$. The last two stages cannot attain this
maximum, by~\eqref{eq:last-core-masses} and~\eqref{eq:supersequence}.
At a maximizing active stage $i$, set $t=m_0-i\ge3$. Then
\[
 \sum_{j=0}^{i+2}S_j
   \le Rn\sum_{s=t-2}^{m_0}u_s
   \le RnP_tu_t=P_tS_i.
\]
Here $S_i>0$, so this is an active fold, a contradiction.
The retained stages are $0,\ldots,m_0-3$, giving the last assertion.
\end{proof}

Write
\[
 T_u=\sum_{s\ge3}u_s,\qquad R_m=\sum_{s\ge m+1}u_s.
\]
For the actual stopped graph $H$ in Lemma~\ref{lem:stopped}, the phase
certificate chooses one finite chain with fee $f_{a,d}$ and integer cap
at least $M_3(a)2^{\theta_2d}$. Thus $H$ has a packing of weight at most
$f_{a,d}n$. Packing the retained cores at cost $1+C\eta$ per unit mass
and adding the retained long cycles gives
\begin{equation}\label{eq:terminal-physical}
 \frac{\wt_\sigma(\mathcal P)}n
       \le f_{a,d}+(1+C\eta)(T_u-R_{m_0})+J_{\rm tr}(a).
\end{equation}
The negative tail $-R_{m_0}$ is kept throughout; it is an unused portion
of a summable majorant, not a negative physical graph cost.

\begin{lemma}[Terminal payment]\label{lem:terminal-payment}
For every phase and every integer $m\ge3$,
\begin{equation}\label{eq:terminal-payment}
 f_{a,d}+(1+C\eta)(T_u-R_m)+J_{\rm tr}(a)
                 \le C\Phi\bigl(a\theta_2^{-(m-3)}\bigr).
\end{equation}
\end{lemma}
\begin{proof}
The phase data first certify
\[
       f_{a,d}+(1+C\eta)T_u+J_{\rm tr}(a)\le C.
\]
For $3\le m<15$, they directly certify
\[
 C L\bigl(a\theta_2^{-(m-3)}\bigr)
    +f_{a,d}+(1+C\eta)T_u+J_{\rm tr}(a)
                      \le C+(1+C\eta)R_m.
\]
These are finite rational interval certificates for the positive series
in Section~\ref{sec:potential}.

For all later $m$, a stronger inequality suffices:
\begin{equation}\label{eq:zero-slack}
            C L\bigl(a\theta_2^{-(m-3)}\bigr)\le(1+C\eta)R_m.
\end{equation}
It is checked separately at $m=15$ and $m=16$. Positivity in
\eqref{eq:tail-recurrence} gives $R_{m+2}\ge\beta R_m$.
On the other hand Lemma~\ref{lem:loss-decay}, applied to every summand of
$L$, gives
\[
 L(x/\theta_2^2)\le
  2^{-(r_*/b)x(\theta_2^{-2}-1)}L(x).
\]
The certificate at $x=a\theta_2^{-12}$ bounds the exponential factor by
$\beta$. The factor only decreases as $x$ increases. Induction separately
on the two parities therefore proves~\eqref{eq:zero-slack} for every
$m\ge15$. Combine it with the total-budget inequality and use
$\Phi(x)=1-L(x)$ on this range to get~\eqref{eq:terminal-payment}.
\end{proof}

The $40$ closed phase intervals cover $[w_A,w_A/\theta_2]$. At a shared
endpoint either adjacent certificate can be used. For the actual phase,
$X=W\theta_2^{-(m_0-3)}\ge a\theta_2^{-(m_0-3)}$.
Equations~\eqref{eq:terminal-physical} and~\eqref{eq:terminal-payment},
and monotonicity of $\Phi$, show that the no-fold branch also has weight
at most $C\phi(n)n$.
